\documentclass[10pt,a4paper]{article}

\usepackage[T1]{fontenc}
\usepackage[utf8]{inputenc}
\usepackage[ngerman]{babel}
\usepackage[a4paper,margin=18mm]{geometry}
\usepackage{array}
\usepackage{booktabs}
\usepackage[table]{xcolor}
\usepackage{lmodern}
\usepackage{microtype}
\usepackage{fancyhdr}

\definecolor{Navy}{HTML}{17324D}
\definecolor{Blue}{HTML}{2F6FA7}
\definecolor{LightBlue}{HTML}{EAF3FA}
\definecolor{Muted}{HTML}{526575}

\pagestyle{fancy}
\setlength{\headheight}{14pt}
\fancyhf{}
\lhead{\textcolor{Muted}{Lehrercodes}}
\rhead{\textcolor{Muted}{Dekodierung}}
\cfoot{\textcolor{Muted}{\thepage}}
\renewcommand{\headrulewidth}{0.4pt}

\setlength{\parindent}{0pt}
\setlength{\parskip}{4pt}
\renewcommand{\arraystretch}{1.12}

\newcommand{\sectiontitle}[1]{%
  \vspace{3pt}%
  {\large\bfseries\color{Navy}#1}\par
  \vspace{2pt}%
}

\begin{document}

{\LARGE\bfseries\color{Navy} Lehrercodes dekodieren}\par
\vspace{2pt}
{\normalsize\color{Muted}Schlüsselkarte für Sicherungsblätter}\par

\vspace{8pt}
\colorbox{LightBlue}{%
  \parbox{\dimexpr\textwidth-2\fboxsep\relax}{%
    \textbf{\color{Navy}Positionsmuster: \texttt{FXXX-KXAX}}\\[3pt]
    1. Zeichen = Fach \quad
    5. Zeichen = Klasse \quad
    7. Zeichen = Aufgabe
  }%
}

\sectiontitle{So wird ein Code gelesen}

Beim Code \texttt{M7QK-4P2X} werden nur drei festgelegte Positionen
dekodiert:

\begin{center}
\begin{tabular}{>{\bfseries}l c l}
\toprule
Position & Zeichen & Bedeutung \\
\midrule
1 & M & Informatik \\
5 & 4 & Klasse 9 \\
7 & 2 & Aufgabe 1 \\
\bottomrule
\end{tabular}
\end{center}

Die übrigen fünf Zeichen sind Prüf- und Tarnzeichen. Sie werden für die
inhaltliche Dekodierung nicht benötigt, gehören aber zwingend zum gültigen
Code.

\vspace{3pt}
\begin{minipage}[t]{0.25\textwidth}
\sectiontitle{Fach}
\centering
\begin{tabular}{l c}
\toprule
Fach & Code \\
\midrule
Informatik & M \\
Physik & R \\
\bottomrule
\end{tabular}
\end{minipage}
\hfill
\begin{minipage}[t]{0.39\textwidth}
\sectiontitle{Klasse}
\centering
\begin{tabular}{c c@{\hspace{9mm}}c c}
\toprule
Klasse & Code & Klasse & Code \\
\midrule
5  & 2 & 9  & 4 \\
6  & 7 & 10 & X \\
7  & K & 11 & D \\
8  & P & 12 & Q \\
\bottomrule
\end{tabular}
\end{minipage}
\hfill
\begin{minipage}[t]{0.29\textwidth}
\sectiontitle{Aufgabe}
\centering
\begin{tabular}{c c}
\toprule
Aufgabe & Code \\
\midrule
1 & 2 \\
2 & P \\
3 & M \\
4 & 7 \\
5 & X \\
\bottomrule
\end{tabular}
\end{minipage}

\sectiontitle{Aktuell verwendete Codes}

\begin{center}
\rowcolors{2}{LightBlue}{white}
\begin{tabular}{l l c}
\toprule
Material & Bedeutung & Lehrercode \\
\midrule
Sicherungsblatt Aufgabe 1 & Informatik, Klasse 9, Aufgabe 1 &
\texttt{M7QK-4P2X} \\
Sicherungsblatt Aufgabe 2 & Informatik, Klasse 9, Aufgabe 2 &
\texttt{MXDT-4QPR} \\
\bottomrule
\end{tabular}
\end{center}

\vfill
\textcolor{Muted}{\small
Diese Schlüsselkarte ist für Lehrkräfte bestimmt. Die Codierung verschleiert
die Bedeutung eines Codes, stellt aber ohne serverseitige Prüfung keine
kryptografische Zugriffssicherung dar.}

\end{document}
